\documentclass[a4paper,12pt]{article}

\usepackage[utf8]{inputenc}
\usepackage[T1]{fontenc}
\usepackage{amsmath}
\usepackage{booktabs}
\usepackage{graphicx}
\usepackage{geometry}
\geometry{margin=2.5cm}
\usepackage{xcolor}
\usepackage{listings}
\usepackage{hyperref}
\hypersetup{colorlinks=true, linkcolor=blue!50!black, urlcolor=blue!50!black}
\renewcommand{\tablename}{Tabela}
\renewcommand{\lstlistingname}{Código}

% ---------- Configuração do realce de sintaxe (C) ----------
\definecolor{codebg}{RGB}{247,247,247}
\definecolor{codekw}{RGB}{0,0,180}
\definecolor{codecomment}{RGB}{0,128,0}
\definecolor{codestring}{RGB}{163,21,21}
\definecolor{framecolor}{RGB}{200,200,200}

\lstdefinestyle{cstyle}{
    language=C,
    backgroundcolor=\color{codebg},
    basicstyle=\ttfamily\footnotesize,
    keywordstyle=\color{codekw}\bfseries,
    commentstyle=\color{codecomment}\itshape,
    stringstyle=\color{codestring},
    numberstyle=\tiny\color{gray},
    numbers=left,
    numbersep=6pt,
    frame=single,
    rulecolor=\color{framecolor},
    breaklines=true,
    breakatwhitespace=false,
    showstringspaces=false,
    tabsize=4,
    captionpos=b,
    columns=flexible,
    morekeywords={size_t},
    extendedchars=true,
    literate=
        {á}{{\'a}}1 {à}{{\`a}}1 {â}{{\^a}}1 {ã}{{\~a}}1
        {é}{{\'e}}1 {ê}{{\^e}}1 {í}{{\'i}}1
        {ó}{{\'o}}1 {ô}{{\^o}}1 {õ}{{\~o}}1
        {ú}{{\'u}}1 {ü}{{\"u}}1 {ç}{{\c c}}1
        {Á}{{\'A}}1 {À}{{\`A}}1 {Â}{{\^A}}1 {Ã}{{\~A}}1
        {É}{{\'E}}1 {Ê}{{\^E}}1 {Í}{{\'I}}1
        {Ó}{{\'O}}1 {Ô}{{\^O}}1 {Õ}{{\~O}}1
        {Ú}{{\'U}}1 {Ü}{{\"U}}1 {Ç}{{\c C}}1
}

\title{Processamento de Lista Encadeada com Tasks em OpenMP}
\author{Israel Soares de Castro Filho \\ Matrícula: 20250070780}
\date{}

\begin{document}

\maketitle

\section{Introdução}

Este relatório descreve a implementação e a análise de um programa em linguagem C que combina uma estrutura de dados clássica -- a lista encadeada -- 
com o modelo de paralelismo baseado em tarefas (\textit{tasks}) da biblioteca OpenMP. 

O objetivo é construir uma lista encadeada em que cada nó armazena o nome de um arquivo fictício e, 
dentro de uma região paralela, percorrer essa lista criando uma tarefa (\texttt{\#pragma omp task}) responsável por processar cada nó individualmente. 

Cada tarefa imprime o nome do arquivo que processou e o identificador da thread que a executou, 
obtido por meio da função \texttt{omp\_get\_thread\_num()}. A partir dos resultados obtidos, 
este trabalho também discute questões relacionadas à correção do processamento paralelo, como a garantia de que cada nó seja processado exatamente uma vez e por uma única tarefa.

\section{Metodologia}

A lista encadeada é construída sequencialmente, antes da entrada na região paralela, 
com dez nós, cada um armazenando o nome de um arquivo fictício no formato \texttt{arquivo\_NN.dat}.

Dentro da região \texttt{\#pragma omp parallel}, o percurso da lista é feito por uma única thread, 
delimitada pela diretiva \texttt{\#pragma omp single nowait}. 

Essa escolha é fundamental: caso todas as threads do time percorressem a lista, cada uma criaria uma tarefa para cada nó, 
resultando em processamento duplicado (o nó seria processado tantas vezes quanto o número de threads). A cláusula \texttt{nowait} evita uma barreira desnecessária 
ao final do \texttt{single}, permitindo que as threads ociosas comecem a executar as tarefas já criadas assim que possível.

Para cada nó visitado, uma tarefa é criada com:

\begin{itemize}
    \item \texttt{\#pragma omp task firstprivate(no\_da\_vez)}. 
\end{itemize}

O uso da cláusula \texttt{firstprivate} é essencial para a correção do programa: ela copia o valor do ponteiro para o nó no exato momento da criação da tarefa. 

Sem essa cláusula, a variável de percurso seria compartilhada por padrão dentro da tarefa e poderia ser alterada pelo laço antes que a tarefa fosse efetivamente executada, 
fazendo com que múltiplas tarefas apontassem para o mesmo nó (tipicamente o último da lista).

Cada tarefa executa a função \texttt{processar\_arquivo}, que imprime o nome do arquivo e o 
identificador da thread corrente (\texttt{omp\_get\_thread\_num()}). Ao final da região paralela, 
a barreira implícita garante que todas as tarefas tenham concluído sua execução antes que a memória da lista seja liberada. 
O código-fonte completo encontra-se no Apêndice~\ref{sec:codigo}.

\section{Resultados}

A execução do programa produziu a seguinte saída, na ordem em que as linhas foram impressas:

\begin{lstlisting}[style=cstyle, language=, numbers=none]
Arquivo 'arquivo_02.dat' processado pela thread 3
Arquivo 'arquivo_05.dat' processado pela thread 1
Arquivo 'arquivo_04.dat' processado pela thread 4
Arquivo 'arquivo_06.dat' processado pela thread 3
Arquivo 'arquivo_03.dat' processado pela thread 5
Arquivo 'arquivo_01.dat' processado pela thread 2
Arquivo 'arquivo_07.dat' processado pela thread 1
Arquivo 'arquivo_08.dat' processado pela thread 2
Arquivo 'arquivo_09.dat' processado pela thread 6
Arquivo 'arquivo_10.dat' processado pela thread 9
\end{lstlisting}

A Tabela~\ref{tab:resultados} reorganiza essa saída por ordem de arquivo, facilitando a verificação de que todos os nós foram processados.

\begin{table}[h]
\centering
\caption{Arquivo processado e thread responsável}
\label{tab:resultados}
\begin{tabular}{@{}cc@{}}
\toprule
\textbf{Arquivo} & \textbf{Thread} \\
\midrule
arquivo\_01.dat & 2 \\
arquivo\_02.dat & 3 \\
arquivo\_03.dat & 5 \\
arquivo\_04.dat & 4 \\
arquivo\_05.dat & 1 \\
arquivo\_06.dat & 3 \\
arquivo\_07.dat & 1 \\
arquivo\_08.dat & 2 \\
arquivo\_09.dat & 6 \\
arquivo\_10.dat & 9 \\
\bottomrule
\end{tabular}
\end{table}

Verifica-se que os dez arquivos (\texttt{arquivo\_01.dat} a \texttt{arquivo\_10.dat}) aparecem exatamente uma vez cada na saída, totalizando as dez linhas esperadas. Nenhum arquivo foi omitido e nenhum foi processado mais de uma vez.

\section{Análise}

\textbf{Cobertura e unicidade do processamento.} Como mostrado na Tabela~\ref{tab:resultados}, cada um dos dez nós da lista foi processado exatamente uma vez, o que confirma a correção da estratégia adotada: a combinação de \texttt{single} (para que apenas uma thread crie as tarefas) com \texttt{firstprivate} (para que cada tarefa capture o ponteiro correto do nó) evita tanto a duplicação quanto a omissão de nós.

\textbf{Distribuição entre threads.} A distribuição das tarefas entre as threads não é uniforme: as threads 1, 2 e 3 processaram dois arquivos cada, enquanto as threads 4, 5, 6 e 9 processaram apenas um arquivo cada. Essa distribuição desigual é esperada, pois o escalonador de tarefas do OpenMP atribui as tarefas dinamicamente às threads ociosas conforme elas terminam seu trabalho anterior -- não há garantia de balanceamento perfeito, especialmente com uma carga de trabalho tão pequena (dez tarefas curtas). A presença de identificadores de thread não sequenciais entre os que aparecem na saída (por exemplo, a ausência das threads 0, 7 e 8 no conjunto observado) apenas reflete que, para essa execução específica, nem todas as threads do time chegaram a receber uma tarefa antes que a lista fosse totalmente percorrida.

\textbf{Ordem de impressão.} A ordem em que as linhas aparecem na saída (por exemplo, \texttt{arquivo\_02.dat} antes de \texttt{arquivo\_01.dat}) não corresponde à ordem dos nós na lista encadeada. 
Isso é uma consequência natural do paralelismo: as tarefas são criadas na ordem da lista, mas sua execução e a impressão do resultado dependem de quando cada thread fica disponível para executá-las. 
Em execuções repetidas do mesmo programa, é esperado que essa ordem varie, assim como pode variar qual thread processa qual arquivo, embora o conjunto de arquivos processados (exatamente uma vez cada) permaneça sempre o mesmo, por construção.

\textbf{Garantia de processamento único.} A robustez do programa depende de três elementos estruturais, e não de sorte na execução: 

\begin{enumerate}
    \item O percurso da lista restrito a uma única thread via \texttt{single}, evitando que múltiplas threads criem tarefas para o mesmo conjunto de nós;
    \item A cláusula \texttt{firstprivate} na variável que referencia o nó atual, garantindo que cada tarefa opere sobre o nó correto mesmo que a variável de laço mude antes da execução efetiva da tarefa;
    \item A barreira implícita ao final da região \texttt{parallel} (que poderia também ser obtida explicitamente com \texttt{\#pragma omp taskwait}), assegurando que todas as tarefas concluam antes que recursos como a lista sejam liberados.
\end{enumerate}

\section{Conclusão}

O experimento confirmou que é possível processar de forma paralela e correta cada nó de uma lista encadeada utilizando tarefas OpenMP, desde que sejam respeitadas as práticas de criação de tarefas em uma região paralela: percurso serial da lista dentro de um bloco \texttt{single} e captura por valor (\texttt{firstprivate}) da referência ao nó em cada tarefa criada. Nos resultados obtidos, todos os dez arquivos foram processados exatamente uma vez, sem duplicações ou omissões, ainda que a ordem de processamento e a thread responsável por cada arquivo variem entre execuções devido à natureza dinâmica do escalonador de tarefas do OpenMP. Esse comportamento é esperado e não compromete a corretude do programa, pois a garantia de unicidade decorre da estrutura do código e não da ordem de execução das threads.

\newpage
\appendix

\section{Código-Fonte}
\label{sec:codigo}

O código-fonte completo utilizado nesta atividade está disponível em \texttt{../src/tarefa7.c} e é reproduzido abaixo.

\lstinputlisting[style=cstyle, caption={Código-fonte completo (\texttt{tarefa7.c})}]{../src/tarefa7.c}

\end{document}